\documentclass[12pt,a4paper]{article}
\usepackage[utf8]{inputenc}
\usepackage[T1]{fontenc}
\usepackage[spanish]{babel}
\usepackage{amsmath}
\usepackage{geometry}
\usepackage{float}
\usepackage{tikz}
\usetikzlibrary{arrows.meta,positioning,shapes.geometric,shadows,fit,backgrounds}
\geometry{left=3cm,right=2.5cm,top=2.5cm,bottom=2.5cm}
\setlength{\parindent}{1.25cm}
\setlength{\parskip}{0.3em}
\definecolor{azulDatos}{RGB}{219,235,255}
\definecolor{azulBorde}{RGB}{42,105,180}
\definecolor{verdeProc}{RGB}{222,246,229}
\definecolor{verdeBorde}{RGB}{40,135,88}
\definecolor{naranjaModelo}{RGB}{255,235,205}
\definecolor{naranjaBorde}{RGB}{211,120,32}
\definecolor{moradoMeta}{RGB}{237,226,255}
\definecolor{moradoBorde}{RGB}{116,80,170}
\definecolor{rojoSalida}{RGB}{255,226,226}
\definecolor{rojoBorde}{RGB}{190,65,65}
\definecolor{grisEval}{RGB}{238,241,245}
\definecolor{grisBorde}{RGB}{95,105,120}
\tikzset{
    bloque/.style={rectangle, rounded corners=5pt, draw=black, align=center, minimum height=1cm, text width=3.2cm, drop shadow},
    dato/.style={bloque, fill=azulDatos, draw=azulBorde, very thick},
    prep/.style={bloque, fill=verdeProc, draw=verdeBorde, very thick},
    modelo/.style={bloque, fill=naranjaModelo, draw=naranjaBorde, very thick},
    meta/.style={bloque, fill=moradoMeta, draw=moradoBorde, very thick},
    salida/.style={bloque, fill=rojoSalida, draw=rojoBorde, very thick},
    eval/.style={bloque, fill=grisEval, draw=grisBorde, very thick},
    proceso/.style={rectangle, rounded corners=4pt, draw=black, align=center, minimum height=1cm, text width=3.4cm, drop shadow},
    hoja/.style={circle, draw=verdeBorde, fill=verdeProc, very thick, align=center, minimum size=1.45cm},
    hojaActiva/.style={circle, draw=rojoBorde, fill=rojoSalida, very thick, align=center, minimum size=1.65cm},
    decision/.style={diamond, draw=black, align=center, aspect=2, text width=2.4cm},
    flecha/.style={-Latex, thick, draw=grisBorde},
    flechaAzul/.style={-Latex, very thick, draw=azulBorde},
    flechaVerde/.style={-Latex, very thick, draw=verdeBorde},
    flechaNaranja/.style={-Latex, very thick, draw=naranjaBorde},
    flechaMorada/.style={-Latex, very thick, draw=moradoBorde},
    flechaRoja/.style={-Latex, very thick, draw=rojoBorde}
}

\begin{document}

\section{Marco teórico}

\subsection{Internet de las Cosas y superficie de ataque en redes IoT}

El Internet de las Cosas (IoT) se comprende como un ecosistema de dispositivos físicos conectados que capturan, procesan e intercambian información mediante redes de comunicación. En este entorno participan sensores, cámaras, televisores inteligentes, equipos portátiles, termostatos, actuadores y otros dispositivos embebidos que interactúan con servicios locales o en la nube. Su valor tecnológico reside en la automatización, el monitoreo remoto y la toma de decisiones basada en datos; sin embargo, estas mismas características incrementan la exposición frente a accesos no autorizados, denegaciones de servicio, reconocimiento de red, explotación de credenciales y control remoto de dispositivos comprometidos (Hindy et al., 2018; Sama et al., 2024).

La seguridad en IoT presenta una complejidad mayor que en redes convencionales debido a tres factores principales. Primero, los dispositivos suelen poseer recursos limitados de procesamiento, memoria y energía, por lo que no siempre pueden ejecutar mecanismos de defensa pesados. Segundo, existe heterogeneidad de fabricantes, protocolos y configuraciones, lo que dificulta establecer políticas uniformes de protección. Tercero, el tráfico IoT puede mostrar patrones repetitivos de operación normal, pero también variaciones abruptas cuando el dispositivo es atacado o integrado a una botnet. Por ello, la detección de intrusiones en IoT requiere modelos capaces de reconocer desviaciones en flujos de red y, especialmente, identificar ataques poco frecuentes que pueden tener alto impacto operativo (Almotairi et al., 2024).

En la presente investigación, las redes IoT se asumen como el dominio de aplicación del problema de clasificación. Cada registro del tráfico de red representa una instancia que puede corresponder a comportamiento benigno o a una categoría de ataque. Esta definición permite trasladar el problema de ciberseguridad a un problema de aprendizaje supervisado, donde el objetivo es asignar correctamente una etiqueta de clase a partir de características extraídas del flujo de comunicación.

\subsection{Sistemas de Detección de Intrusiones en redes IoT}

Un Sistema de Detección de Intrusiones (IDS) es un mecanismo orientado a identificar actividades anómalas o maliciosas dentro de una red o sistema informático. En redes IoT, los IDS pueden analizar paquetes, flujos, registros de eventos o patrones de comportamiento de los dispositivos. Desde el punto de vista de la técnica de detección, suelen clasificarse en enfoques basados en firmas, anomalías o aprendizaje automático (Hindy et al., 2018; Diwan et al., 2021).

Los IDS basados en firmas comparan el tráfico observado con patrones conocidos de ataque. Aunque son útiles frente a amenazas previamente registradas, presentan limitaciones ante ataques nuevos, variantes de malware o comportamientos que no coinciden exactamente con una regla definida. Los IDS basados en anomalías, en cambio, modelan el comportamiento normal y reportan desviaciones significativas. Esta estrategia resulta pertinente en IoT porque muchos ataques modifican la frecuencia, duración, tamaño o dirección de los flujos de red. No obstante, su desempeño depende de la calidad de los datos y de la capacidad del modelo para diferenciar entre variación legítima y actividad maliciosa (Adewole et al., 2025).

El aprendizaje automático permite construir IDS supervisados a partir de conjuntos de datos etiquetados. En este enfoque, los algoritmos aprenden relaciones entre variables predictoras y clases de tráfico. Para la presente tesis, el IDS se concibe como un modelo de clasificación multiclase o binario, según el nivel de agrupación de los ataques, entrenado con registros del dataset CICIoT2023. Su eficacia se evalúa mediante métricas de clasificación y su viabilidad mediante indicadores computacionales.

\subsection{Ataques minoritarios y desbalance de clases}

El desbalance de clases ocurre cuando una o varias categorías del conjunto de datos poseen una cantidad de instancias considerablemente menor que otras. En ciberseguridad, este fenómeno es frecuente porque el tráfico benigno o ciertos ataques masivos suelen dominar la distribución, mientras que ataques específicos, raros o de baja frecuencia aparecen con menor representación. El problema no es únicamente estadístico, sino operativo: una clase minoritaria puede corresponder a una amenaza crítica, aunque su frecuencia sea reducida (Shanmugam et al., 2024).

En modelos de clasificación, el desbalance puede inducir sesgo hacia la clase mayoritaria. Un algoritmo puede alcanzar una exactitud global alta clasificando correctamente la mayoría de registros frecuentes, pero fallar en la detección de ataques escasos. En un IDS, este comportamiento es riesgoso porque eleva los falsos negativos; es decir, permite que ataques reales sean clasificados como tráfico benigno o como una clase menos crítica. Por esta razón, la presente investigación prioriza métricas sensibles al desempeño de las clases minoritarias, especialmente Recall y F1-Score, antes que depender únicamente de la Accuracy (Shanmugam et al., 2024; Wu et al., 2025).

El índice de desbalance permite cuantificar la relación entre clases mayoritarias y minoritarias. Si $N_{max}$ representa el número de instancias de la clase mayoritaria y $N_{min}$ el número de instancias de una clase minoritaria, el índice puede expresarse como:

\[
IR = \frac{N_{max}}{N_{min}}
\]

Un valor elevado de $IR$ indica mayor asimetría en la distribución. En el contexto de esta tesis, reducir dicho índice mediante técnicas de rebalanceo contribuye a que los modelos base reciban mayor información sobre la frontera de decisión de las clases menos representadas.

\subsection{Dataset CICIoT2023 como base empírica del estudio}

El CICIoT2023 es un conjunto de datos moderno orientado a la investigación en seguridad IoT, detección de intrusiones y análisis de anomalías. Fue construido a partir de un banco de pruebas que simula comunicaciones de dispositivos IoT reales y escenarios de ataque contemporáneos. Incluye tráfico benigno y múltiples familias de amenazas, como ataques de denegación de servicio, denegación de servicio distribuida, fuerza bruta, botnets, reconocimiento y otras intrusiones asociadas a redes de dispositivos conectados. Su uso como base empírica resulta consistente con antecedentes nacionales que lo emplean para clasificación de anomalías en redes IoT (Raico Gallardo, 2023).

Para esta investigación, el CICIoT2023 cumple una función metodológica central porque proporciona registros etiquetados y características numéricas de flujo de red adecuadas para aprendizaje supervisado. Las variables disponibles describen propiedades como duración del flujo, tamaño de paquetes, tasas de transmisión, banderas, conteos y medidas estadísticas derivadas. Estas características permiten entrenar modelos que no dependen del contenido semántico del paquete, sino del comportamiento observable del flujo, lo cual resulta pertinente para escenarios IoT.

El conjunto de datos empleado en la investigación fue dividido en tres subconjuntos: entrenamiento con $5\,491\,971$ registros y 47 variables, validación con $1\,176\,851$ registros y 47 variables, y prueba con $1\,176\,851$ registros y 47 variables. Esta partición permite separar el ajuste del modelo, la selección de configuraciones y la evaluación final. De este modo, se reduce el riesgo de sobreajuste y se favorece una medición más objetiva del desempeño del modelo híbrido propuesto.

\subsection{Preprocesamiento y representación de características de flujo}

El preprocesamiento de datos es una etapa crítica para construir modelos de detección de intrusiones. En conjuntos masivos como CICIoT2023, los registros pueden contener valores atípicos, duplicados, variables con escalas heterogéneas o características con bajo aporte predictivo. Por ello, antes del entrenamiento se requiere limpiar, transformar y preparar los datos para que los algoritmos de clasificación procesen patrones consistentes.

Las características de flujo representan estadísticamente la comunicación entre dispositivos. A diferencia del análisis profundo de paquetes, que puede ser costoso y sensible a cifrado, el análisis por flujo resume el comportamiento de la conexión mediante atributos agregados. Esta representación favorece la escalabilidad porque permite procesar grandes volúmenes de tráfico y, al mismo tiempo, conserva señales relevantes para distinguir tráfico benigno de tráfico malicioso.

En la presente tesis, el tratamiento de características se vincula directamente con la viabilidad del modelo. Al trabajar con millones de registros, la arquitectura debe mantener equilibrio entre rendimiento predictivo y costo computacional. Por ello, el marco teórico adopta una perspectiva de clasificación basada en flujos, aprendizaje supervisado y evaluación experimental sobre datos previamente particionados.

\subsection{Técnicas de rebalanceo y Borderline-SMOTE}

Las técnicas de rebalanceo buscan modificar la distribución de clases para mejorar el aprendizaje sobre categorías poco representadas. Existen estrategias de submuestreo, que reducen instancias de la clase mayoritaria; sobremuestreo, que incrementa instancias de la clase minoritaria; y métodos híbridos que combinan ambas aproximaciones. En detección de intrusiones, el rebalanceo debe aplicarse con cuidado, ya que eliminar demasiados registros mayoritarios puede descartar información útil y generar pérdida de representatividad (Musthafa et al., 2024; Nigar et al., 2025).

SMOTE (Synthetic Minority Oversampling Technique) genera instancias sintéticas de la clase minoritaria mediante interpolación entre ejemplos cercanos. En lugar de duplicar registros, produce nuevos puntos dentro del espacio de características. Borderline-SMOTE es una variante que se concentra en las instancias minoritarias ubicadas cerca de la frontera de decisión, es decir, aquellas que tienen mayor probabilidad de ser confundidas con clases mayoritarias. Esta característica lo vuelve pertinente para ataques minoritarios, porque refuerza las zonas donde el clasificador suele cometer errores (Sayegh et al., 2024).

Formalmente, si $x_i$ es una instancia minoritaria y $x_{nn}$ uno de sus vecinos más cercanos de la misma clase, una instancia sintética puede generarse como:

\[
x_{new} = x_i + \lambda (x_{nn} - x_i)
\]

donde $\lambda$ es un valor aleatorio en el intervalo $[0,1]$. En Borderline-SMOTE, la selección de $x_i$ prioriza muestras minoritarias situadas en regiones de frontera. En la investigación, esta técnica se relaciona con la primera dimensión de la variable independiente: balanceo de datos. Sus indicadores son el índice de desbalance y la distribución de datos en la frontera de decisión.

\subsection{Aprendizaje automático supervisado para detección de intrusiones}

El aprendizaje automático supervisado utiliza datos etiquetados para estimar una función que relacione variables de entrada con una salida conocida. En detección de intrusiones, esta salida corresponde a la clase del tráfico: benigno o ataque, o bien una familia específica de ataque. La función aprendida se evalúa posteriormente sobre datos no vistos para estimar su capacidad de generalización.

Los modelos supervisados aplicados a IDS deben responder a dos exigencias. La primera es predictiva: reconocer patrones complejos y diferenciar clases similares. La segunda es operacional: ejecutar entrenamiento e inferencia en tiempos razonables. Esta doble exigencia justifica el uso de algoritmos basados en árboles y métodos de ensamble, porque suelen manejar variables tabulares de alta dimensión, capturar relaciones no lineales y ofrecer buen rendimiento en clasificación (Sama et al., 2024; Asim, 2024).

En este estudio, el aprendizaje supervisado no se limita a entrenar un clasificador único. Se adopta una arquitectura híbrida en la que modelos base generan predicciones y probabilidades, y un meta-modelo aprende a combinar dichas salidas. Esta decisión teórica y arquitectónica responde al problema central: mejorar la detección de ataques minoritarios reduciendo errores residuales de los modelos individuales.

\subsection{Árboles de decisión y modelos basados en ensamble}

Los árboles de decisión son modelos predictivos que dividen recursivamente el espacio de características mediante reglas. Cada nodo interno representa una condición sobre una variable, cada rama representa el resultado de la condición y cada hoja corresponde a una predicción. Su principal ventaja es la interpretabilidad, ya que permite traducir decisiones en reglas comprensibles.

Los modelos de ensamble combinan múltiples clasificadores para obtener un predictor más robusto. La lógica central es que la combinación de varios modelos puede reducir errores que aparecerían en un clasificador individual. Entre los enfoques más usados se encuentran Bagging, Boosting y Stacking. Bagging reduce varianza mediante entrenamiento paralelo sobre muestras del conjunto de datos. Boosting entrena modelos secuenciales que corrigen errores de iteraciones previas. Stacking combina modelos de primer nivel mediante un meta-modelo que aprende a integrar sus predicciones (Lin et al., 2025; Elasaad et al., 2023).

La presente investigación se posiciona en el enfoque de Stacking porque el problema no consiste solo en obtener alta exactitud global, sino en aprovechar fortalezas complementarias de XGBoost y LightGBM para mejorar el reconocimiento de clases minoritarias. Además, el uso de un Árbol de Decisión como meta-modelo de Nivel 1 aporta reglas de combinación más interpretables que una caja negra profunda.

\subsection{XGBoost}

XGBoost (Extreme Gradient Boosting) es un algoritmo de aprendizaje conjunto basado en árboles de decisión y optimización por gradiente. Su funcionamiento consiste en construir árboles de manera secuencial, donde cada nuevo árbol intenta corregir los errores cometidos por los anteriores. Esta estrategia permite modelar relaciones no lineales y capturar interacciones complejas entre variables.

En clasificación, XGBoost minimiza una función objetivo compuesta por una función de pérdida y un término de regularización:

\[
Obj = \sum_{i=1}^{n} l(y_i,\hat{y}_i) + \sum_{k=1}^{K} \Omega(f_k)
\]

donde $l(y_i,\hat{y}_i)$ mide el error entre la etiqueta real y la predicción, mientras que $\Omega(f_k)$ penaliza la complejidad de cada árbol. Esta regularización ayuda a controlar el sobreajuste, aspecto importante cuando se trabaja con millones de registros y múltiples clases de ataque.

XGBoost resulta adecuado para la investigación porque suele obtener alto rendimiento en datos tabulares, admite ponderación de clases, maneja variables con relaciones complejas y ofrece mecanismos de control de profundidad, tasa de aprendizaje y número de estimadores. En el modelo propuesto, XGBoost actúa como clasificador base de Nivel 0 y aporta probabilidades de clase que posteriormente son utilizadas por el meta-modelo. Su selección es coherente con estudios comparativos que reportan rendimiento competitivo de modelos de boosting en IDS para IoT (Sama et al., 2024).

\begin{figure}[H]
\centering
\resizebox{0.98\textwidth}{!}{%
\begin{tikzpicture}[node distance=1.05cm and 1.2cm]
\node[dato] (data) {CICIoT2023\\entrenamiento\\rebalanceado};
\node[prep, right=of data] (pred0) {Predicción inicial\\$\hat{y}^{(0)}$};
\node[modelo, right=of pred0] (grad) {Gradientes\\errores residuales};
\node[modelo, below=of grad] (tree) {Nuevo árbol\\$f_t(x)$};
\node[meta, left=of tree] (reg) {Regularización\\$\Omega(f_t)$};
\node[modelo, left=of reg] (update) {Actualización\\$\hat{y}^{(t)}$};
\node[salida, below=of update] (prob) {Probabilidades\\por clase};
\node[eval, right=of prob] (stack) {Entrada al\\Stacking};

\draw[flechaAzul] (data) -- (pred0);
\draw[flechaVerde] (pred0) -- (grad);
\draw[flechaNaranja] (grad) -- (tree);
\draw[flechaMorada] (tree) -- (reg);
\draw[flechaNaranja] (reg) -- (update);
\draw[flechaRoja] (update) -- (prob);
\draw[flechaMorada] (prob) -- (stack);
\draw[flechaNaranja] (update.north) |- node[pos=0.26, above]{corrige errores} (grad.south);

\begin{scope}[on background layer]
\node[draw=naranjaBorde, fill=naranjaModelo, rounded corners=8pt, opacity=0.18, fit=(grad)(tree)(reg)(update), inner sep=8pt] {};
\end{scope}
\end{tikzpicture}
}
\caption{Funcionamiento general de XGBoost como clasificador base del modelo híbrido.}
\end{figure}

\noindent\textbf{Pseudocódigo de XGBoost aplicado al proyecto}

\begin{quote}
\small
\begin{enumerate}
    \item Ingresar el conjunto de entrenamiento de CICIoT2023 después de limpieza y Borderline-SMOTE.
    \item Inicializar la predicción $\hat{y}^{(0)}$ para cada flujo de red.
    \item Para cada iteración $t=1,\ldots,T$, calcular gradientes y errores de clasificación.
    \item Entrenar un árbol $f_t(x)$ que corrija los errores residuales de la iteración anterior.
    \item Penalizar árboles demasiado complejos mediante el término de regularización $\Omega(f_t)$.
    \item Actualizar la predicción con tasa de aprendizaje $\eta$.
    \item Generar probabilidades de clase para cada registro, priorizando la detección de ataques minoritarios.
    \item Enviar dichas probabilidades al nivel de Stacking como metacaracterísticas.
\end{enumerate}
\end{quote}

\subsection{LightGBM}

LightGBM (Light Gradient Boosting Machine) es un algoritmo de boosting basado en árboles diseñado para mejorar eficiencia y escalabilidad. A diferencia de métodos que crecen árboles por nivel, LightGBM utiliza crecimiento por hoja, seleccionando la hoja con mayor ganancia para expandirla. Esta estrategia puede lograr menor pérdida con menos divisiones, aunque requiere controlar parámetros como profundidad máxima, número de hojas y tasa de aprendizaje para evitar sobreajuste.

Su pertinencia en esta investigación se sustenta en el tamaño del dataset CICIoT2023. Al trabajar con más de siete millones de registros distribuidos entre entrenamiento, validación y prueba, el algoritmo debe ser capaz de procesar datos masivos con tiempos competitivos. LightGBM es especialmente útil cuando se requiere reducir tiempo de entrenamiento y latencia de inferencia sin sacrificar capacidad predictiva (Ammari, 2024; Almotairi et al., 2024).

En la arquitectura planteada, LightGBM complementa a XGBoost en el Nivel 0. Mientras XGBoost aporta robustez y regularización fuerte, LightGBM aporta velocidad y escalabilidad. La combinación de ambos permite que el meta-modelo reciba señales predictivas diversas, lo cual fortalece el enfoque de Stacking.

\begin{figure}[H]
\centering
\resizebox{0.88\textwidth}{!}{%
\begin{tikzpicture}[node distance=1.2cm and 1.15cm]
\node[dato] (hist) {Histogramas\\de variables};
\node[modelo, right=of hist] (gain) {Cálculo de\\ganancia};
\node[hojaActiva, below=of gain] (root) {Raíz\\máx.\\ganancia};
\node[hoja, below left=1.1cm and 1.8cm of root] (a) {Hoja A\\estable};
\node[hojaActiva, below right=1.1cm and 1.8cm of root] (b) {Hoja B\\mayor\\pérdida};
\node[hoja, below left=1.1cm and 0.7cm of b] (b1) {B1};
\node[hojaActiva, below right=1.1cm and 0.7cm of b] (b2) {B2\\candidata};
\node[salida, right=1.45cm of b2] (out) {Árbol ligero\\probabilidades\\rápidas};

\draw[flechaAzul] (hist) -- (gain);
\draw[flechaNaranja] (gain) -- (root);
\draw[flechaVerde] (root) -- (a);
\draw[flechaRoja] (root) -- node[pos=0.55, right]{se expande} (b);
\draw[flechaVerde] (b) -- (b1);
\draw[flechaRoja] (b) -- node[pos=0.55, right]{siguiente} (b2);
\draw[flechaMorada] (b2) -- (out);

\begin{scope}[on background layer]
\node[draw=verdeBorde, fill=verdeProc, rounded corners=8pt, opacity=0.16, fit=(root)(a)(b)(b1)(b2), inner sep=10pt] {};
\end{scope}
\end{tikzpicture}
}
\caption{Crecimiento por hoja en LightGBM para reducir pérdida con menor costo computacional.}
\end{figure}

\noindent\textbf{Pseudocódigo de LightGBM aplicado al proyecto}

\begin{quote}
\small
\begin{enumerate}
    \item Ingresar los registros de entrenamiento de CICIoT2023 con sus 47 variables.
    \item Construir histogramas o agrupaciones internas de valores para acelerar el cálculo de divisiones.
    \item Calcular la ganancia de cada posible división en las hojas disponibles.
    \item Seleccionar la hoja con mayor ganancia y expandirla.
    \item Controlar la complejidad mediante número máximo de hojas, profundidad y tasa de aprendizaje.
    \item Repetir el proceso hasta alcanzar el número de árboles o el criterio de parada.
    \item Producir probabilidades de clase con baja latencia de inferencia.
    \item Entregar las probabilidades al meta-modelo de Stacking junto con las de XGBoost.
\end{enumerate}
\end{quote}

\subsection{Generalización por apilamiento o Stacking}

Stacking, o generalización por apilamiento, es una arquitectura de ensamble que combina las predicciones de varios modelos base mediante un meta-modelo. En un esquema de dos niveles, los clasificadores de Nivel 0 aprenden directamente de las variables originales, mientras que el clasificador de Nivel 1 aprende de las salidas generadas por los modelos base. Estas salidas pueden incluir etiquetas predichas, probabilidades de clase o medidas derivadas de incertidumbre. Este patrón arquitectónico se alinea con la tendencia de IDS híbridos y de múltiples etapas reportada en literatura reciente (Popoola et al., 2025; Elasaad et al., 2023).

La lógica del Stacking es que distintos algoritmos pueden cometer errores diferentes. Si el meta-modelo aprende cuándo confiar más en un clasificador u otro, puede mejorar la decisión final. Para evitar fuga de información, las predicciones usadas para entrenar el meta-modelo deben generarse de manera controlada, preferentemente mediante validación cruzada o particiones separadas.

En esta tesis, la arquitectura se define de la siguiente manera: en el Nivel 0 se entrenan XGBoost y LightGBM sobre el conjunto de entrenamiento rebalanceado; luego, ambos generan probabilidades de clase sobre datos de validación; finalmente, en el Nivel 1, un Árbol de Decisión utiliza estas probabilidades y metadatos de incertidumbre, como la entropía de Shannon, para producir la clasificación final. Esta estructura es congruente con la variable independiente de la matriz de consistencia: modelo híbrido basado en Stacking, XGBoost y LightGBM.

\begin{figure}[H]
\centering
\resizebox{0.98\textwidth}{!}{%
\begin{tikzpicture}[node distance=1.0cm and 1.15cm]
\node[dato] (raw) {CICIoT2023\\Train\\Validación\\Test};
\node[prep, right=of raw] (clean) {Limpieza\\codificación\\normalización};
\node[prep, right=of clean] (smote) {Borderline-\\SMOTE\\clases raras};

\node[modelo, below left=1.35cm and 0.55cm of smote] (xgb) {Nivel 0\\XGBoost\\robustez};
\node[modelo, below right=1.35cm and 0.55cm of smote] (lgbm) {Nivel 0\\LightGBM\\rapidez};
\node[meta, below=2.55cm of smote] (features) {Meta-\\características\\$P_{XGB}$, $P_{LGBM}$\\entropía $H$};
\node[meta, below=of features] (dt) {Nivel 1\\Árbol de\\Decisión};
\node[salida, below=of dt] (final) {Salida final\\benigno o\\tipo de ataque};
\node[eval, right=of final] (metrics) {Validación\\Recall\\F1-Score\\latencia};

\draw[flechaAzul] (raw) -- (clean);
\draw[flechaVerde] (clean) -- (smote);
\draw[flechaNaranja] (smote) -- (xgb);
\draw[flechaNaranja] (smote) -- (lgbm);
\draw[flechaMorada] (xgb) -- node[pos=0.45,left]{$P_{XGB}$} (features);
\draw[flechaMorada] (lgbm) -- node[pos=0.45,right]{$P_{LGBM}$} (features);
\draw[flechaMorada] (features) -- (dt);
\draw[flechaRoja] (dt) -- (final);
\draw[flecha] (final) -- (metrics);

\begin{scope}[on background layer]
\node[draw=verdeBorde, fill=verdeProc, rounded corners=8pt, opacity=0.12, fit=(clean)(smote), inner sep=8pt] {};
\node[draw=naranjaBorde, fill=naranjaModelo, rounded corners=8pt, opacity=0.14, fit=(xgb)(lgbm), inner sep=8pt] {};
\node[draw=moradoBorde, fill=moradoMeta, rounded corners=8pt, opacity=0.15, fit=(features)(dt), inner sep=8pt] {};
\end{scope}
\end{tikzpicture}
}
\caption{Arquitectura de Stacking propuesta para la detección de ataques minoritarios en redes IoT.}
\end{figure}

\noindent\textbf{Pseudocódigo del enfoque de Stacking propuesto}

\begin{quote}
\small
\begin{enumerate}
    \item Dividir CICIoT2023 en entrenamiento, validación y prueba.
    \item Aplicar limpieza, transformación de variables y Borderline-SMOTE sobre el conjunto de entrenamiento.
    \item Entrenar XGBoost y LightGBM como modelos base de Nivel 0.
    \item Obtener probabilidades de clase de ambos modelos sobre el conjunto de validación.
    \item Calcular la entropía de Shannon de las probabilidades para representar incertidumbre predictiva.
    \item Construir una matriz de metacaracterísticas con probabilidades de XGBoost, probabilidades de LightGBM y entropía.
    \item Entrenar el Árbol de Decisión de Nivel 1 usando la matriz de metacaracterísticas.
    \item Evaluar el modelo final sobre el conjunto de prueba mediante Recall, F1-Score, Accuracy, tiempo de entrenamiento, latencia de inferencia y memoria RAM.
\end{enumerate}
\end{quote}

\subsection{Probabilidad de clase e incertidumbre predictiva}

En clasificación probabilística, un modelo no solo asigna una etiqueta, sino una distribución de probabilidad sobre las posibles clases. Si existen $C$ clases, la predicción para una instancia puede representarse como:

\[
P(y|x) = [p_1, p_2, \ldots, p_C]
\]

donde $p_c$ indica la probabilidad estimada de que la instancia pertenezca a la clase $c$ y se cumple que $\sum_{c=1}^{C} p_c = 1$. Estas probabilidades son útiles para el Stacking porque contienen más información que una etiqueta final. Por ejemplo, dos modelos pueden predecir la misma clase, pero con niveles de confianza distintos.

La incertidumbre predictiva mide cuán segura o dispersa es la distribución de probabilidades. Una predicción con probabilidad dominante cercana a uno indica alta confianza; una distribución distribuida entre varias clases indica ambigüedad. En ataques minoritarios, esta información es relevante porque muchas instancias se ubican cerca de fronteras de decisión y pueden confundirse con tráfico benigno o con ataques mayoritarios (Sepúlveda-Fontaine y Amigó, 2024).

\subsection{Entropía de Shannon como metadato de incertidumbre}

La entropía de Shannon es una medida de incertidumbre asociada a una distribución de probabilidad. En clasificación, permite cuantificar el grado de dispersión de las probabilidades asignadas por un modelo. Su uso en análisis de datos y aprendizaje automático se fundamenta en la teoría de la información, donde la entropía permite representar incertidumbre, variabilidad y contenido informacional de una distribución (Sepúlveda-Fontaine y Amigó, 2024). Para una instancia con $C$ clases, la entropía se expresa como:

\[
H(X) = -\sum_{c=1}^{C} p_c \log_2(p_c)
\]

Cuando una clase concentra casi toda la probabilidad, la entropía es baja; cuando las probabilidades se distribuyen entre varias clases, la entropía es alta. Por ello, la entropía funciona como un metadato útil para el meta-modelo de Stacking. No reemplaza la probabilidad de clase, sino que resume el nivel de duda de los clasificadores base.

En esta investigación, la entropía de Shannon forma parte de los indicadores de la dimensión arquitectura de apilamiento. Su incorporación permite que el Árbol de Decisión de Nivel 1 no solo reciba las probabilidades generadas por XGBoost y LightGBM, sino también una señal explícita sobre la incertidumbre de dichas predicciones. Esto es especialmente pertinente para ataques minoritarios, donde las regiones de frontera suelen producir distribuciones probabilísticas menos definidas.

\subsection{Métricas de evaluación de la detección}

La evaluación de un IDS requiere métricas que reflejen la capacidad real de detección. La matriz de confusión organiza los resultados en verdaderos positivos, falsos positivos, verdaderos negativos y falsos negativos. En el contexto del presente proyecto, estos valores se interpretan sobre la clasificación de tráfico IoT del dataset CICIoT2023, priorizando las clases de ataque minoritarias. Así, $TP$ (verdaderos positivos) representa las intrusiones minoritarias correctamente identificadas por el modelo híbrido; $FN$ (falsos negativos) contabiliza los ataques reales que el sistema omite o clasifica erróneamente como tráfico benigno u otra clase no correspondiente; $FP$ (falsos positivos) agrupa registros benignos o de otra clase que fueron marcados incorrectamente como ataque minoritario; y $TN$ (verdaderos negativos) corresponde a los registros no pertenecientes a la clase evaluada que fueron descartados correctamente. A partir de estos valores se calculan indicadores como Accuracy, Recall, Precision y F1-Score.

La Accuracy mide la proporción total de clasificaciones correctas:

\[
Accuracy = \frac{TP + TN}{TP + TN + FP + FN}
\]

Aunque es útil como indicador global, puede ser engañosa en conjuntos desbalanceados. Por ello, la presente investigación prioriza el Recall y el F1-Score para evaluar ataques minoritarios. El Recall mide la proporción de intrusiones reales correctamente detectadas por la arquitectura basada en Stacking, XGBoost y LightGBM:

\[
Recall = \frac{TP}{TP + FN}
\]

Donde $TP$ verifica las intrusiones minoritarias clasificadas correctamente por el modelo, mientras que $FN$ contabiliza las omisiones de ataque, es decir, los flujos maliciosos que no fueron detectados como tales. Maximizar matemáticamente esta función valida la capacidad del modelo híbrido para reducir falsos negativos, aspecto central en la detección de ataques de baja frecuencia en redes IoT. El F1-Score combina Precision y Recall mediante media armónica:

\[
F1 = 2 \cdot \frac{Precision \cdot Recall}{Precision + Recall}
\]

En esta ecuación, $Precision$ expresa la proporción de predicciones de ataque que realmente corresponden a intrusiones, mientras que $Recall$ expresa la proporción de ataques reales detectados. Maximizar el F1-Score permite validar el equilibrio entre detectar la mayor cantidad de ataques minoritarios y evitar alertas incorrectas. En la tesis, Recall, F1-Score y Accuracy operacionalizan la dimensión eficacia de clasificación de la variable dependiente.

\subsection{Eficiencia computacional: tiempo de entrenamiento y latencia de inferencia}

La eficacia predictiva de un IDS debe complementarse con su eficiencia computacional. En redes IoT, un modelo con alta precisión pero excesiva latencia puede ser inviable, porque la detección tardía reduce la capacidad de respuesta ante intrusiones. Por ello, el desempeño computacional se define como una dimensión de evaluación de la variable dependiente (Almotairi et al., 2024; Popoola et al., 2025).

El tiempo de entrenamiento corresponde al periodo requerido para ajustar los parámetros del modelo sobre los datos de entrenamiento. Este indicador es importante en la etapa experimental y en escenarios donde el IDS debe actualizarse periódicamente con nuevos datos. La latencia de inferencia, en cambio, mide el tiempo requerido para clasificar una instancia o lote de instancias nuevas. Este indicador se vincula directamente con la posibilidad de uso operativo del modelo.

En la investigación, la comparación entre el modelo híbrido y los modelos base Random Forest y XGBoost permitirá determinar si la mejora en Recall y F1-Score mantiene tiempos de entrenamiento e inferencia competitivos. De este modo, el modelo propuesto no será evaluado únicamente por exactitud, sino por su equilibrio entre detección y viabilidad de despliegue.

\subsection{Validación estadística mediante Wilcoxon}

La prueba de rangos con signo de Wilcoxon es una prueba no paramétrica utilizada para comparar dos conjuntos de mediciones relacionadas cuando no se asume normalidad en las diferencias. En el contexto de esta tesis, puede utilizarse para contrastar el rendimiento del modelo híbrido frente a modelos base sobre métricas como Recall, F1-Score o Accuracy obtenidas en distintas particiones, clases o repeticiones experimentales.

La hipótesis nula establece que no existe diferencia significativa entre el desempeño del modelo propuesto y el modelo de comparación. La hipótesis alternativa plantea que sí existe una diferencia significativa. Si el valor $p$ obtenido es menor que el nivel de significancia definido, se rechaza la hipótesis nula y se concluye que la diferencia observada no se debe únicamente al azar.

Esta prueba se integra al marco teórico porque conecta la evaluación empírica con la hipótesis general de la investigación: que el modelo híbrido basado en Stacking, XGBoost y LightGBM mejora significativamente la detección de ataques minoritarios en redes IoT.

\subsection{Modelo teórico propuesto}

El modelo teórico de la investigación articula cuatro componentes: datos IoT, rebalanceo, clasificación híbrida y evaluación. Primero, el dataset CICIoT2023 proporciona registros de tráfico benigno y malicioso con 47 variables por instancia. Segundo, Borderline-SMOTE actúa sobre el conjunto de entrenamiento para mejorar la representación de clases minoritarias ubicadas en la frontera de decisión. Tercero, XGBoost y LightGBM aprenden patrones de clasificación como modelos base de Nivel 0. Cuarto, un Árbol de Decisión de Nivel 1 integra probabilidades de clase y entropía de Shannon para emitir la predicción final.

La relación entre variables puede expresarse de forma conceptual así:

\[
Y = f(M_{stacking}(XGB, LGBM, H, P), BSMOTE, X)
\]

donde $Y$ representa la detección de ataques minoritarios, $X$ los registros de tráfico de red, $BSMOTE$ el proceso de rebalanceo, $P$ las probabilidades de clase generadas por los modelos base y $H$ la entropía de Shannon como medida de incertidumbre. Esta formulación sintetiza el posicionamiento arquitectónico del estudio: la mejora de detección no depende de un único algoritmo, sino de la integración entre balanceo de datos, clasificadores de boosting y metadatos probabilísticos dentro de una arquitectura de Stacking.

\subsection{Operacionalización teórica de las variables}

La variable independiente es el modelo híbrido basado en Stacking, XGBoost y LightGBM. Conceptualmente, se define como una arquitectura de aprendizaje supervisado que combina clasificadores de boosting en un primer nivel y un meta-modelo de decisión en un segundo nivel. Operacionalmente, se mide mediante dos dimensiones: arquitectura de apilamiento y técnica de balanceo de datos. La primera incluye como indicadores la probabilidad de clase, la entropía de Shannon y los hiperparámetros del meta-modelo. La segunda incluye el índice de desbalance y la distribución de datos en la frontera de decisión tras aplicar Borderline-SMOTE.

La variable dependiente es la detección de ataques minoritarios en redes IoT. Conceptualmente, se entiende como la capacidad del sistema para identificar correctamente clases de ataque de baja frecuencia dentro del tráfico de red. Operacionalmente, se mide mediante dos dimensiones: eficacia de clasificación y desempeño computacional. La eficacia de clasificación se evalúa mediante Recall, F1-Score y Accuracy. El desempeño computacional se evalúa mediante tiempo de entrenamiento, latencia de inferencia y consumo de memoria RAM.

Esta operacionalización asegura la correspondencia entre el marco teórico, la matriz de consistencia y la rúbrica técnica. Los conceptos no se presentan como definiciones aisladas, sino como componentes medibles del modelo experimental que será implementado sobre CICIoT2023.

\end{document}
